\appendix
\section{Data Snapshot and Counting Protocol}
\label{sec:appendix-protocol}

To keep counts reproducible as the live corpora evolve, we separate
paper numbers (frozen cohort for this draft) from live recounts.
The counting protocol is versioned in the repository:
\texttt{docs/corpus\_protocol.md}.

\paragraph{Counting tool.}
Session counting is performed by
\texttt{tools/corpus\_counts.py}, which classifies Claude Code JSONL
files into \texttt{main}, \texttt{subagent}, \texttt{compact},
\texttt{prompt\_suggestion}, and \texttt{other}, and reports both
raw and content-hash-deduplicated totals.

\paragraph{Recount command.}
\begin{verbatim}
./tools/reproduce_paper_counts.sh
\end{verbatim}
This writes a timestamped snapshot JSON to
\texttt{paper/data/corpus\_snapshot\_YYYYMMDD\_live.json}.

\paragraph{5-minute reproduction path.}
From repository root:
\begin{verbatim}
./tools/reproduce_paper_counts.sh
python3 tools/check_paper_numbers.py \
  --paper paper/main.tex \
  --snapshot paper/data/corpus_snapshot_20260307_live.json
latexmk -pdf paper/main.tex
\end{verbatim}
Expected outputs: (1) a new snapshot JSON under \texttt{paper/data/},
(2) a paper-number consistency report, and (3) \texttt{paper/main.pdf}.
On a typical laptop this path completes in minutes; the recount step is
I/O-bound on local corpus size.

\paragraph{Current live recount artifact (2026-03-07).}
\texttt{paper/data/corpus\_snapshot\_20260307\_live.json}
with roots:
\texttt{\~/.claude/projects} and
\texttt{\~/projects/yanantin/tmp/ubuntu-vm.claude/projects},
size filter \texttt{min\_size=10000}, and Claude Desktop summary
from \texttt{\~/projects/yanantin/tmp/claude-desktop/conversations.json}.

For this paper version, headline results use a frozen cohort to
avoid drift across revisions; live recount artifacts are provided for
auditability.

% ====================================================================
